%==============================================================
% BAB IV  HASIL DAN PEMBAHASAN
%==============================================================

\chapter{HASIL DAN PEMBAHASAN}

Dalam penulisan hasil dan pembahasan dapat dipisah sebagai bab Hasil dan bab
Pembahasan, atau digabung menjadi bab Hasil dan Pembahasan.
Pemisahan atau penggabungan kedua bab ini bergantung pada bidang studi,
atau sesuai dengan arahan pembimbing.

% Jika dipisah, judul bab diganti:
%   \chapter{HASIL}
% dan file baru:
%   \chapter{PEMBAHASAN}


%--------------------------------------------------------------
\section{Hasil}
%--------------------------------------------------------------

Sajikan hasil penelitian secara sistematis.
Tabel dan gambar digunakan untuk mempermudah pembacaan data.
Setiap tabel dan gambar dirujuk dalam teks sebelum atau sesudah kemunculannya,
misalnya: (Tabel~\ref{tab:pisang_ambon}) atau (Gambar~\ref{fig:contoh_gambar1}).

Penulisan satuan mengikuti sistem SI.
Contoh penulisan satuan dalam teks: kekerasan buah diukur dalam
mm 50 g$^{-1}$ detik$^{-1}$ \citep{Kochel2005}.


%--------------------------------------------------------------
\section{Pembahasan}
%--------------------------------------------------------------

Pembahasan mengaitkan hasil penelitian dengan teori, hipotesis, dan hasil
penelitian sebelumnya.
Argumentasi penulis harus jelas dan didukung oleh data.
Kebaruan atau temuan penting ditekankan di bagian ini \citep{Onlamoon2010}.

Jangan hanya mengulang hasil; jelaskan mengapa hasil tersebut terjadi dan
apa artinya dalam konteks bidang ilmu.
Gunakan pustaka yang relevan untuk mendukung interpretasi.
